\section{Poisson normalization}
\label{sec:normalization}

The prescribed first-shell derivative determines the remaining scalar
normalization.  Poisson summation first equates the values at the origin;
its differentiated dilation identity then proves that their common value is
positive, without estimating that value numerically.

Let the exact coefficient lists be those of
Proposition~\ref{interpolation:solution}, and put
\[
 G_i(s)=e^{-\pi hs}H_i(s),\qquad F(x)=G_1(b|x|^2).
\]
Section~\ref{sec:fourier} proves that $F$ is real, radial, and Schwartz,
with $\widehat F(\xi)=G_2(b|\xi|^2)$.  The interpolation conditions give
\begin{equation}\label{normalization:jets}
 \begin{split}
 G_1(n)&=G_2(n)=0,\\
 G_2'(n)&=0,\qquad
 G_1'(n)=-Q_1e^{-\pi h}\mathbf1_{\{n=1\}}
 \qquad(n\in\mathcal N).
 \end{split}
\end{equation}
Indeed, $G_i'(n)=e^{-\pi hn}H_i'(n)$ because $H_i(n)=0$.

\subsection{The dilation identity}

Use the lattice $\Lambda$ of covolume one and the shell coordinate $s_a$ from
\eqref{fourier:lattice}--\eqref{fourier:shell-coordinate}.  Every positive
$s_a$ lies in $\mathcal N$.  Exactly six lattice vectors have
$s_a=1$: from
$4(j^2+jk+k^2)=(2j+k)^2+3k^2=4$, one obtains
\[
 (j,k)=(1,0),\ (-1,0),\ (0,1),\ (0,-1),\ (1,-1),\ (-1,1).
\]

As shown in Section~\ref{sec:fourier}, $\Lambda^*=\mathcal R\Lambda$.
Apply Poisson summation on $\Lambda$ to
$F_t(x)=F(\sqrt t\,x)$, whose Fourier transform is
$t^{-1}\widehat F(\xi/\sqrt t)$.  Radiality makes the sums over
$\Lambda^*$ and $\Lambda$ equal, and gives, for every $t>0$,
\begin{equation}\label{normalization:poisson}
 \sum_{a\in\Lambda}G_1(ts_a)
 =\frac1t\sum_{a\in\Lambda}G_2(s_a/t).
\end{equation}

We will differentiate this identity twice.  Here is the needed convergence
justification.  On a compact subinterval of $(0,\infty)$, each of the first
two $t$ derivatives of $F(\sqrt t\,x)$ is a finite sum of spatial derivatives
of $F$ multiplied by polynomials in $x$ and bounded functions of $t$.
The same statement holds for $t^{-1}\widehat F(x/\sqrt t)$.
The Schwartz bounds therefore imply, for every $M$, a common bound
$C_M(1+|x|)^{-M}$ for these functions and their first two $t$ derivatives.
The number of points of a fixed planar lattice in a radius-$R$ disk is
$O(R^2)$, so $\sum_{a\in\Lambda}(1+|a|)^{-M}$ converges for $M>2$.
Uniform convergence on compact $t$ intervals now permits both termwise
differentiations of \eqref{normalization:poisson}.

\begin{lemma}[Exact origin value]\label{normalization:origin-lemma}
The Fourier pair constructed above satisfies
\begin{equation}\label{normalization:origin}
 F(0)=\widehat F(0)=6Q_1e^{-\pi h}>0.
\end{equation}
\end{lemma}

\begin{proof}
At $t=1$, all nonzero summands in \eqref{normalization:poisson} vanish by
\eqref{normalization:jets}.  Hence $G_1(0)=G_2(0)$.
The derivative of its left side at $1$ is
\[
 \sum_{a\in\Lambda}s_aG_1'(s_a)=-6Q_1e^{-\pi h},
\]
because only the six vectors on the first shell contribute.  On the right,
\[
 \left.\frac{\dd}{\dd t}\left(t^{-1}G_2(s/t)\right)\right|_{t=1}
 =-G_2(s)-sG_2'(s).
\]
For $s=s_a>0$, both terms vanish by the exact jets, while $s_a=0$ contributes
$-G_2(0)$.  Equality of the derivatives gives
$G_2(0)=6Q_1e^{-\pi h}$.  The positivity follows from $Q_1>0$ in
\eqref{fourier:node-bounds}.  This calculation uses the Fourier pairing and
the exact interpolation conditions, but not the global sign estimates.
\end{proof}

\subsection{The final dilation}

\begin{proof}[Proof of Theorem~\ref{intro:main}]
Define
\begin{equation}\label{normalization:scaling}
 f(x)=\frac{F(x/\sqrt b)}{bF(0)},\qquad
 \widehat f(\xi)=\frac{\widehat F(\sqrt b\,\xi)}{F(0)}.
\end{equation}
The second identity is the two-dimensional Fourier scaling rule.  Since
$F(0)>0$, the dilation and normalization preserve real-valuedness,
radiality, and Schwartz regularity.  They give
$\widehat f(0)=1$ and $f(0)=1/b=2/\sqrt3$.
By Proposition~\ref{signs:global} and the positive factor $e^{-\pi hs}$,
$G_2(s)=e^{-\pi hs}H_2(s)\ge0$ for $s\ge0$ and
$G_1(s)=e^{-\pi hs}H_1(s)\le0$ for $s\ge1$.
Thus $\widehat f\ge0$ everywhere.  If $|x|\ge1$,
then $b|x/\sqrt b|^2=|x|^2\ge1$, so $f(x)\le0$.  These are exactly the
four assertions of \eqref{intro:conditions}.
\end{proof}

The scaling also records the prescribed contact data in the final radial
coordinates.  Write $f_{\mathrm{rad}}(r)=f(r,0)$ and
$\widehat f_{\mathrm{rad}}(\rho)=\widehat f(\rho,0)$ for $r,\rho\ge0$.  Then
\[
 f_{\mathrm{rad}}(r)=\frac{G_1(r^2)}{bF(0)},\qquad
 \widehat f_{\mathrm{rad}}(\rho)=\frac{G_2(b^2\rho^2)}{F(0)}.
\]
Proposition~\ref{signs:global} therefore gives the exact zero sets in the
respective sign domains:
\begin{equation}\label{normalization:zero-sets}
 \begin{aligned}
 \{r^2:r\ge1,\ f_{\mathrm{rad}}(r)=0\}&=\mathcal N,\\
 \{\rho^2:\rho\ge0,\ \widehat f_{\mathrm{rad}}(\rho)=0\}
   &=b^{-2}\mathcal N=\frac43\mathcal N.
 \end{aligned}
\end{equation}
The physical zero at $r=1$ is simple, every other listed physical zero is
exactly double, and every listed Fourier zero is exactly double.  Indeed,
the Gaussian and scalar factors do not vanish, while the changes of
variable $s=r^2$ and $s=b^2\rho^2$ have nonzero derivatives at these positive
radii, so they preserve the orders established for $H_1,H_2$.
At the exclusion radius, the radial derivative is
\begin{equation}\label{normalization:radial-slope}
 f_{\mathrm{rad}}'(1)
 =\frac{2G_1'(1)}{bF(0)}=-\frac{1}{3b}=-\frac{2}{3\sqrt3}.
\end{equation}

\subsection{The packing consequence}

The relation between the signs and packing density is particularly visible
for a periodic packing.  Suppose the center set is
$\bigcup_{j=1}^N(L+t_j)$, where $L$ has covolume $V$, the $t_j$ are distinct
modulo $L$, and distinct centers have distance at least one.  Absolute
convergence follows from the Schwartz bounds.  Every summand below except
the $N$ diagonal terms at the origin is nonpositive, so
\[
 \mathcal S:=\sum_{j,k=1}^N\sum_{v\in L}f(v+t_j-t_k)\le Nf(0).
\]
Poisson summation on each translated lattice gives
\[
 \mathcal S=\frac1V\sum_{w\in L^*}\widehat f(w)
       \left|\sum_{j=1}^N e^{2\pi i w\cdot t_j}\right|^2
 \ge\frac{N^2}{V}\widehat f(0),
\]
because every Fourier summand is nonnegative.  It follows that
$N/V\le f(0)/\widehat f(0)=2/\sqrt3$.  The disks have area $\pi/4$, hence
their covered-area density is at most $\pi/(2\sqrt3)$.

For an arbitrary, not necessarily periodic, packing, the same implication
is Theorem~3.1 of Cohn and Elkies~\cite{CohnElkies2003}.  Its admissibility
hypotheses hold because both $f$ and $\widehat f$ are Schwartz.
For disks of radius $1/2$, the Cohn--Elkies center-density normalization is
one quarter of the center-number density in their upper-density convention.
Multiplying it by the unit-disk area $\pi$ therefore gives the corresponding
covered-area density, equivalently $\pi/4$ times that center-number density.
Their upper density takes a supremum over observation centers, so it bounds
the fixed-origin limsup used in the introduction.
For the fixed-origin convention and separated centers, the
covered area inside $B_R$ lies between $(\pi/4)$ times the number of centers
in $B_{R-1/2}$ and $(\pi/4)$ times the number in $B_{R+1/2}$; division by
$\pi R^2$ and passage to the limsup matches the convention in the
introduction.  The lattice generated by $(1,0)$ and $(1/2,b)$ has one center
per area $b$ and attains this value.  This is the classical packing-density
theorem recovered from the new Fourier certificate, not a new density claim.

The restriction of the physical zero set to integer squared radii also
recovers the classical periodic equality case.  We use the even integral
lattice argument of Cohn and Elkies~\cite[Section~8]{CohnElkies2003}.

\begin{corollary}[Periodic equality case]\label{normalization:periodic-uniqueness}
If a periodic packing of disks of radius $1/2$ has covered-area density
$\pi/(2\sqrt3)$, then its center set is the image of
$\ZZ(1,0)+\ZZ(1/2,b)$ under a Euclidean isometry.
\end{corollary}

\begin{proof}
Write the center set as
$\mathcal C=\bigcup_{j=1}^N(L+t_j)$, with the notation of the periodic
packing argument above.  Equality of its density with $\pi/(2\sqrt3)$ gives
$N/V=f(0)/\widehat f(0)$.  Thus both inequalities in the Poisson comparison
are equalities.  Absolute convergence permits us to write the defect in
the physical inequality as
\[
 0=Nf(0)-\mathcal S
 =\sum_{\substack{1\le j,k\le N,\ v\in L\\
                   v\ne0\ \text{or}\ j\ne k}}
       -f(v+t_j-t_k).
\]
Every summand is nonnegative, so every one is zero.  Any difference of
distinct centers occurs in this sum, and \eqref{normalization:zero-sets}
therefore gives
\[
 |x-y|^2\in\mathcal N\subset\ZZ
 \qquad(x,y\in\mathcal C,\ x\ne y).
\]
Only the physical zero set is used here.

Translate one center to the origin and dilate by $\sqrt2$.  Put
\[
 L_0=\sqrt2\,L,\qquad
 \mathcal C_0=\sqrt2\,(\mathcal C-t_1),\qquad
 \Gamma=\operatorname{span}_{\ZZ}\mathcal C_0.
\]
Here $\Gamma$ consists of finite integer linear combinations of the
scaled centers, and $L_0\subset\mathcal C_0$ because the first coset becomes
$L_0$.  All squared distances in $\mathcal C_0$ are even integers.  Since
$0\in\mathcal C_0$, polarization gives
\[
 \langle x,y\rangle
 =\frac{|x|^2+|y|^2-|x-y|^2}{2}\in\ZZ
 \qquad(x,y\in\mathcal C_0).
\]
Integer linear combinations preserve these integral pairings, and their
squared norms remain even because the cross terms carry a factor of two.
Thus every pairing in $\Gamma$ is integral and every squared norm in
$\Gamma$ is even.

In particular, pairing with a basis of $L_0$ shows
$L_0\subset\Gamma\subset L_0^*$.  Both $L_0$ and $L_0^*$ are lattices of
rank two, so $L_0^*/L_0$ is finite.  Hence $\Gamma$ is a finite union of
$L_0$ cosets and is itself a lattice of rank two.  Put
$d:=[\Gamma:L_0]$.  A basis of $\Gamma$ has a Gram matrix
\[
 \begin{pmatrix}a&m\\m&c\end{pmatrix},
 \qquad a,c\in2\ZZ_{>0},\quad m\in\ZZ.
\]
Its determinant $ac-m^2=\operatorname{covol}(\Gamma)^2$ is a positive
integer congruent to $0$ or $3$ modulo $4$.  It is therefore at least $3$,
and $\operatorname{covol}(\Gamma)\ge\sqrt3$.
On the other hand, $\mathcal C_0$ consists of $N$ distinct $L_0$ cosets,
all among the $d$ cosets of $\Gamma$.  Since
$\operatorname{covol}(L_0)=2V$, we obtain
\[
 \sqrt3\le\operatorname{covol}(\Gamma)
 =\frac{2V}{d}\le\frac{2V}{N}=\sqrt3.
\]
Equality holds throughout, so $d=N$ and the two sets of cosets coincide:
$\mathcal C_0=\Gamma$.

It remains to identify this lattice of Gram determinant $3$.  Choose a
shortest nonzero vector $w_1\in\Gamma$ and extend it to a lattice basis
$w_1,w_2$; this is possible because a shortest vector is primitive.
Put $a=|w_1|^2$, and replace $w_2$ by an integer translate along $w_1$
so that $m=\langle w_1,w_2\rangle$ satisfies $|m|\le a/2$.
With $c=|w_2|^2\ge a$ and the even integer $a\ge2$,
\[
 3=ac-m^2\ge\frac34a^2.
\]
Hence $a=2$.  The equation $2c-m^2=3$, with $c$ even and $|m|\le1$,
forces $c=2$ and $m=\pm1$.  This is the Gram matrix of the triangular
lattice with shortest length $\sqrt2$.  Scaling back by $1/\sqrt2$ proves
the assertion.
\end{proof}

\subsection{A second-derivative consistency identity}

The same dilation identity contains an additional relation among the exact
curvatures.  It is not needed for the construction or its sign proof.
For $n\in\mathcal N$, let
$r_\Lambda(n)=\#\{a\in\Lambda:s_a=n\}$, allowing the value zero for nodes
that are not lattice radii; in particular, $r_\Lambda(1)=6$.

\begin{proposition}[Second differentiated Poisson identity]
\label{normalization:second-poisson}
The exact pair satisfies the absolutely convergent identity
\begin{equation}\label{normalization:curvature-identity}
 \sum_{n\in\mathcal N}r_\Lambda(n)n^2
       \bigl(G_1''(n)-G_2''(n)\bigr)=2G_2(0).
\end{equation}
Using the global signs, it follows that
\begin{equation}\label{normalization:curvature-inequality}
 H_1''(1)-H_2''(1)\ge2(1-\pi h)Q_1.
\end{equation}
\end{proposition}

\begin{proof}
The convergence argument above permits a second differentiation of
\eqref{normalization:poisson}.  At $t=1$,
\[
 \left.\frac{\dd^2}{\dd t^2}\left(t^{-1}G_2(s/t)\right)\right|_{t=1}
 =2G_2(s)+4sG_2'(s)+s^2G_2''(s).
\]
On the left the corresponding term is $s^2G_1''(s)$.  At every positive
node, the first two terms on the right vanish by
\eqref{normalization:jets}; at the origin they contribute $2G_2(0)$.
Grouping the remaining absolutely convergent lattice sums by their squared
radii proves \eqref{normalization:curvature-identity}.

For a positive node $n>1$, $G_1$ has a local maximum equal to zero and
$G_2$ has a local minimum equal to zero, by Proposition~\ref{signs:global}
and $e^{-\pi hs}>0$.
Therefore $G_1''(n)\le0\le G_2''(n)$.  All terms in
\eqref{normalization:curvature-identity} beyond the first shell are
nonpositive.  Since $r_\Lambda(1)=6$ and $G_2(0)=6Q_1e^{-\pi h}$,
\[
 G_1''(1)-G_2''(1)\ge2Q_1e^{-\pi h}.
\]
Finally, differentiating $G_i=e^{-\pi hs}H_i$ twice and inserting the
exact value and slope at $1$ gives
\[
 G_1''(1)-G_2''(1)
 =e^{-\pi h}\bigl(H_1''(1)-H_2''(1)+2\pi hQ_1\bigr).
\]
This yields \eqref{normalization:curvature-inequality}.  Unlike the origin
calculation, this last inequality is conditional on the global signs that
were proved in Section~\ref{sec:signs}.
\end{proof}
